\documentclass[10pt,a4paper]{amsart}
\usepackage[margin=19mm]{geometry}
\usepackage{amsmath,amssymb,mathtools}
\usepackage[scr=rsfs,cal=euler]{mathalfa}
\usepackage{xfrac}
\usepackage[unicode,hidelinks,bookmarks=false]{hyperref}
\newcommand{\ie}{i.e.\ }
\newcommand{\cf}{cf.\ }
\newcommand{\resp}{respectively\ }
\newcommand{\Subsection}{\subsection}
\providecommand{\nobreakdash}{\nobreak\hbox{-}}
\setlength{\emergencystretch}{2em}
\multlinegap0pt
\usepackage{bm}
\usepackage{bbm}
\usepackage{dsfont}
\usepackage{xspace}
\usepackage{cancel}
\usepackage{mathtools}
\usepackage{amsthm}
\makeatletter
\@ifundefined{theorem}{\newtheorem{theorem}{Theorem}}{}
\@ifundefined{dcases}{\newtheorem{dcases}{Dcases}}{}
\@ifundefined{definition}{\newtheorem{definition}[theorem]{Definition}}{}
\@ifundefined{lemma}{\newtheorem{lemma}[theorem]{Lemma}}{}
\makeatother
\def\N{\mathbb{N}}
\def\R{\mathbb{R}}
\def\loc{{\rm loc}}
\def\sym{{\rm sym}}
\newcommand{\E}{\mathcal{E}}
\newcommand{\st}{ \ : \ }
\newcommand{\vphi}{\varphi}
\title[Functions of bounded Musielak-Orlicz-type deformation and an]{Functions of bounded Musielak-Orlicz-type deformation and anisotropic Total Generalized Variation for image-denoising problems}
\author{Giacomo Bertazzoni, Elisa Davoli, Samuele Ricc\`o, Elvira Zappale}
\date{Source: arXiv:2509.09237v4 (CC BY 4.0)}
\begin{document}
\maketitle

\noindent\emph{Source and licence.} Functions of bounded Musielak-Orlicz-type deformation and anisotropic Total Generalized Variation for image-denoising problems. Version arXiv:2509.09237v4 (CC BY 4.0), distributed under the Creative Commons Attribution 4.0 International licence, as recorded in the arXiv licence metadata for this exact version and shown on the abstract page https://arxiv.org/abs/2509.09237v4. Full rights notice: licenses/arxiv-2509.09237-CC-BY-4.0.txt.\par\smallskip

\noindent\textit{Editorial scope.} Counted unit: the Lemma whose source label is \texttt{lem:BV-gradient} in the article (source0.tex lines 848--858, source label \texttt{lem:BV-gradient}), delivered together with the article's own complete original proof of it (source0.tex lines 859--892, one \texttt{proof} environment of 34 lines). No mathematics is written, repaired or re-derived: statement and proof are the article's own text line for line. Required context retained uncounted: def:dual (lines 466--472). Every definition, display and label of those lines is retained unchanged. Omitted from this package and not redistributed: 1--465, 473--847, 893--1983. Nothing inside a retained excerpt is omitted or altered: every retained body is delivered exactly as the article's lines are written, so no omission receipt of any kind accompanies this package. Citations: the retained excerpts carry no citation command at all, so no bibliography entry of the article is transcribed and no reference is redistributed. Reference adaptation: none was needed and none was made; every reference inside the delivered unit resolves inside it, so no unresolved reference or citation is produced. Printed slips kept literal: no printed word, symbol, hypothesis or number of the counted statement or of its proof was altered. Layout and class: the article's own class is \texttt{article} with packages including times, lmodern, fontenc, inputenc, babel, amssymb, amsmath, mathtools, amsthm, amsfonts, mathrsfs, latexsym, color, graphicx; the wrapper is the lane's pinned \texttt{amsart} editorial wrapper at 10pt on A4, which restates the theorem-like environments the retained text needs and adds the article's own macro definitions that the retained text needs (\texttt{\textbackslash E}, \texttt{\textbackslash N}, \texttt{\textbackslash R}, \texttt{\textbackslash loc}, \texttt{\textbackslash st}, \texttt{\textbackslash sym}, \texttt{\textbackslash vphi}), with extra wrapper packages (bm, bbm, dsfont, xspace, cancel, mathtools). The delivered numbering is the unit's own. Assets: no retained span contains a figure environment, a graphics-inclusion command, a PDF-page inclusion, a table or a TikZ picture; no image file is copied into this package, the package declares no asset dependency and no image file is redistributed with it. Licence: version arXiv:2509.09237v4 of this article is distributed under the Creative Commons Attribution 4.0 International licence, as recorded in the arXiv licence metadata for this exact version and shown on the abstract page https://arxiv.org/abs/2509.09237v4. A full rights notice is delivered at licenses/arxiv-2509.09237-CC-BY-4.0.txt. Characters: every retained excerpt is pure ASCII, so the delivered engine, XeLaTeX with its default Latin Modern text font, reports no missing character.
\begin{definition}
\label{def:dual}
    Let $\vphi \in \Phi_w(\Omega)$. We define the associate space $(L^\vphi)'(\Omega)=\{u \in L^0(\Omega):  \|u\|_{(L^\vphi)'}<+\infty\}$ with
    \begin{equation*}
        \|u\|_{(L^\vphi)'} := \sup_{\|v\|_\vphi \le 1} \int_\Omega uv \, dx.
    \end{equation*}
\end{definition}

\begin{lemma}
\label{lem:BV-gradient}
    Let $\Omega \subset \R^n$ be a bounded, connected and open set. Let $\vphi \in \Phi_w(\Omega)$ and let $u \in LD_\loc(\Omega)$. Then
    \begin{equation*}
        \widetilde{V}^{n \times n}_\vphi(Eu) \le \| \E u \|_{(L^{\vphi^*}(\Omega; \R^{n \times n}))'}.
    \end{equation*}
    Moreover, if additionally $C^1_c(\Omega; \R^{n \times n})$ is dense in $L^{\vphi^*}(\Omega; \R^{n \times n})$, then
    \begin{equation*}
        \widetilde{V}^{n \times n}_\vphi(Eu) = \| \E u \|_{(L^{\vphi^*}(\Omega;\R^{n \times n}))'}.
    \end{equation*}
\end{lemma}

\begin{proof}
    Since $u \in LD_\loc(\Omega)$, it follows from the definition of $\widetilde{V}^{n \times n}_\vphi$ that 
    \begin{equation}
    \label{eq:integrationByParts}
        \widetilde{V}^{n \times n}_\vphi(Eu) = \sup \left\{ \int_\Omega \E u : w \, dx \st w \in C^1_c(\Omega; \R^{n \times n}_\sym), \|w\|_{\vphi^*} \le 1 \right\}. 
    \end{equation}
    The definition of the associate space norm (extended in a standard way to the vector valued setting), see Definition \ref{def:dual}, implies that, chosen $w \in C^1_c(\Omega; \R^{n \times n}_\sym)$ such that $\|w\|_{\vphi^*} \le 1$,
    \begin{equation*}
        \int_\Omega \E u : w \, dx \le \int_\Omega |\E u| |w| \, dx \le \| \E u\|_{(L^{\vphi^*}(\Omega;\R^{n \times n}))'} \|w\|_{\vphi^*}.
    \end{equation*}
    Taking the supremum over those $w$, we conclude that $\widetilde{V}^{n \times n}_\vphi(\E u) \le \|\E u\|_{(L^{\vphi^*}(\Omega;\R^{n \times n}))'}$. 
    \\
    Under the density assumption, namely $C^1_c(\Omega;\R^{n \times n})$ being dense in $L^{\vphi^*}(\Omega;\R^{n \times n})$, we next show the opposite inequality, namely $ \|\E u\|_{(L^{\vphi^*}(\Omega;\R^{n \times n}))'} \le \widetilde{V}^{n \times n}_\vphi(\E u)$. Indeed, let $w \in L^{\vphi^*}(\Omega; \R^{n \times n}_\sym)$ with $\|w\|_{\vphi^*} = 1$ and let $(w_j)_j \subset C^1_c(\Omega; \R^{n \times n}_\sym)$ be a sequence such that $w_j \to w$ in $L^{\vphi^*}(\Omega; \R^{n \times n}_\sym)$ and pointwise almost everywhere. Since we know that also $w_j/\|w_j\|_{\vphi^*} \to w$ in $L^{\vphi^*}(\Omega; \R^{n \times n}_\sym)$, we may assume that $\|w_j\|_{\vphi^*} = 1$ for every $j \in \N$. By Fatou's Lemma, we have
    \begin{equation*}
        \liminf_{j \to \infty} \int_\Omega \E u : w_j \, dx \ge \int_\Omega \E u : w \, dx.
    \end{equation*}
    So, from \eqref{eq:integrationByParts} follows that 
    \begin{equation*}
        \widetilde{V}^{n \times n}_\vphi(\E u) \ge \sup \left\{ \int_\Omega \E u : w \, dx \st w \in L^{\vphi^*}(\Omega; \R^{n \times n}_\sym), \| w\|_{\vphi^*} \le 1 \right\}.
    \end{equation*}
    Let $h\in L^{\vphi^*}(\Omega)$. We can set
    \begin{equation*}
        w :=
        \begin{dcases}
            \frac{\mathcal E u}{|\mathcal E u|} \, h & \textnormal{ if } |\mathcal E u|\neq 0, \\
            0 & \textnormal{ otherwise.}
        \end{dcases}
    \end{equation*}
    This gives 
    \begin{equation*}
        \widetilde{V}^{n \times n}_\vphi( \E u) \ge \sup \left\{ \int_\Omega |\E u| \, h \, dx \st h \in L^{\vphi^*}(\Omega), \| h\|_{\vphi^*}\le 1 \right\} = \| \E u \|_{(L^{\vphi^*}(\Omega;\R^{n \times n}))'}.
    \end{equation*}
    Hence $\widetilde{V}^{n \times n}_\vphi(\E u) = \| \E u \|_{(L^{\vphi^*}(\Omega;\R^{n \times n}))'}$. This concludes the proof.
\end{proof}

\end{document}
